Multimodal XR systems can observe task activity, while wearable sensing and proactive assistants can provide information about individual users and their assistance needs ~\cite{p37,pro2,guide,sensing}. Interpersonal physiology also offers evidence about collaborative processes, but its meaning depends on the task and cannot be reduced to a universal synchrony score ~\cite{p32,p49,huang}. We propose a physiology-informed agent that combines signals from several participants with speech, gaze, interaction traces, and shared task events to decide whether, whom, and how to assist. The central question is whether relational physiological information improves these decisions beyond behavioral and task information alone. The proposed evaluation separates support-need detection, intervention selection, system feasibility, and effects on collaboration. General XR collaboration provides the primary setting; medical team training and collaborative education define extensions with distinct communication and learning objectives. In the extensions, unassisted transfer is required to distinguish immediate assistance from retained capability ~\cite{p62}. The intended contributions are an operational model of team support needs, a reproducible sensing-to-action pipeline, and evidence about the benefits and limits of adding physiology.

\missing{[TO COMPLETE after studies: primary setting, number of teams, system performance, intervention effects with uncertainty, and one evidence-supported conclusion.]}
